\documentclass[11pt]{article}
\usepackage{amsmath}
\usepackage{amsfonts}
\usepackage{amssymb}
\usepackage{amsthm}

\newtheorem{theorem}{Theorem}

\title{An Integer Boundary Condition for Ratios of Strictly Positive Linear Inequalities}
\author{\textbf{Your Full Name Goes Here} \\ Independent Researcher}
\date{\today}

\begin{document}

\maketitle

\begin{abstract}
While the subtraction of simultaneous algebraic equations preserves equality, the subtraction of linear inequalities does not generally preserve order. This paper identifies a specific boundary rule under which a modified property of order is strictly maintained. We prove that for any set of strictly positive integers satisfying given linear order conditions, a division ratio condition mathematically guarantees a distinct minimum bound for the difference of the primary terms.
\end{abstract}

\section{Introduction}
In elementary algebra, it is a well-established principle that given two inequalities $x > y$ and $a > b$, one cannot safely perform the subtraction operation $x - a > y - b$. Because subtraction shifts positions by non-uniform magnitudes rather than scaling proportionally, the directional sign of the resulting statement is vulnerable to inversion based on the arbitrary magnitude of the subtracted terms. This paper establishes a strict boundary condition where a clear inequality property is preserved, provided the domain is bounded to the set of strictly positive integers ($\mathbb{Z}^+$).

\section{The Core Theorem}
\begin{theorem}
Let $a, b, c, d \in \mathbb{Z}^+$ be strictly positive integers such that $a > b$ and $c > d$. If the ratio condition $\frac{a}{c} \ge \frac{b}{d}$ is satisfied, then the discrete difference $a - b$ is bounded below by the opposing ratio:
\begin{equation}
a - b \ge \frac{b}{d}
\end{equation}
\end{theorem}

\begin{proof}
We begin with the primary ratio constraint:
\begin{equation}
\frac{a}{c} \ge \frac{b}{d}
\end{equation}
Because all variables are strictly positive integers, we cross-multiply across the inequality without altering its direction, yielding:
\begin{equation}
a \cdot d \ge b \cdot c
\end{equation}
We are also given that $c > d$. Since $c$ and $d$ are integers, this strict inequality implies:
\begin{equation}
c \ge d + 1
\end{equation}
Multiplying both sides of this inequality by the positive integer $b$ gives $b \cdot c \ge b(d + 1)$. Substituting this lower bound into our cross-multiplication result allows us to chain the inequalities using the \texttt{align*} environment:
\begin{align*}
a \cdot d &\ge b(d + 1) \\
a \cdot d &\ge b \cdot d + b
\end{align*}
To isolate the original fractional relationship, we divide all terms by the positive integer $d$:
\begin{equation}
a \ge b + \frac{b}{d}
\end{equation}
Subtracting $b$ from both sides yields our final desired lower bound:
\begin{equation}
a - b \ge \frac{b}{d}
\end{equation}
This completes the proof.
\end{proof}

\section{Conclusion}
By restricting the domain to positive integers, we demonstrate that while direct inequality subtraction remains chaotic, the cross-ratio condition acts as a strict mathematical filter. This ratio effectively constrains the value of $\frac{b}{d}$ to stay beneath or equal to the discrete integer difference $a - b$, offering a predictable boundary shortcut in discrete number theory.

\end{document}
